% Part: first-order-logic
% Chapter: syntax-and-semantics
% Section: Structures

\documentclass[../../../include/open-logic-section]{subfiles}

\begin{document}

\olfileid{fol}{syn}{str}

\olsection{\printtoken{P}{structure} voor talen van de eerste orde
(eerste-ordetalen)}

\begin{explain}
Op zichzelf hebben de tekens van een taal van de eerste orde nog geen
vastgelegde betekenis. Zo'n taal heet dan \emph{ongeïnterpreteerd}. De
!!{constant}s, !!{function}s en !!{predicate}s wijzen dus nog niets
bepaalds aan. Een betekenis krijgen ze pas als we
\article{structure} \emph{!!{structure}} vastleggen. De structuur
bepaalt eerst het \emph{domein}. Dat zijn de objecten die de
!!{constant}s aanwijzen, waarop de !!{function}s werken en waarover de
kwantoren gaan. Daarna legt de structuur vast welke !!{constant} welk
object aanwijst, hoe elke !!{function} objecten aan objecten koppelt en
voor welke objecten de !!{predicate}s gelden. Met !!^{structure}s
definiëren logici \emph{begrippen over betekenis}, zoals semantisch gevolg,
geldigheid en vervulbaarheid. In boeken en artikelen heten zulke
gegevens niet altijd hetzelfde: men spreekt van ``structuren'',
``interpretaties'' of ``modellen''.
\end{explain}

\begin{defn}[!!^{structure}s]
\Article{structure} \emph{!!{structure}}~$\Struct M$ voor een taal van
de eerste orde $\Lang{L}$ bestaat uit het volgende:
\begin{enumerate}
\item \emph{Domein:} een verzameling $\Domain M$ die niet leeg is.
\item \emph{Wat de !!{constant}s aanwijzen (hun interpretatie):} voor
  iedere !!{constant}~$c$ van $\Lang{L}$ !!a{element}
  $\Assign{c}{M} \in \Domain M$.
\item \emph{Voor welke objecten de !!{predicate}s gelden (hun
  interpretatie):} voor ieder $n$-plaatsig !!{predicate}~$R$ van
  $\Lang{L}$ (behalve $\eq$) een $n$-plaatsige relatie
  $\Assign{R}{M} \subseteq \Domain{M}^n$.
\item \emph{Welke bewerkingen de !!{function}s voorstellen (hun
  interpretatie):} voor ieder $n$-plaatsig !!{function}~$f$ van
  $\Lang{L}$ een $n$-plaatsige functie
  $\Assign{f}{M} \colon \Domain{M}^n \to \Domain{M}$.
\end{enumerate}
\end{defn}

\begin{ex}
!!^a{structure}~$\Struct M$ voor de taal van de rekenkunde moet de
volgende gegevens bevatten. Er is een verzameling als domein. Uit
$\Domain M$ wordt één element $\Assign{\Obj 0}{M}$ gekoppeld aan de
!!{constant}~$\Obj 0$. Er is een eenplaatsige functie
$\Assign{\Obj \prime}{M} \colon \Domain{M} \to \Domain M$. Verder zijn
er twee tweeplaatsige functies, $\Assign{\Obj +}{M}$ en $\Assign{\Obj
\times}{M}$, beide $\Domain M^2 \to \Domain M$, en ten slotte een
tweeplaatsige relatie
$\Assign{\Obj <}{M} \subseteq \Domain{M}^2$.

Een voor de hand liggende keuze is:
\begin{enumerate}
\item $\Domain N = \Nat$
\item $\Assign{\Obj 0}{N} = 0$
\item $\Assign{\Obj \prime}{N}(n) = n + 1$ voor alle $n \in \Nat$
\item $\Assign{\Obj +}{N}(n, m) = n + m$ voor alle $n, m \in \Nat$
\item $\Assign{\Obj \times}{N}(n, m) = n\cdot m$ voor alle $n, m \in \Nat$
\item $\Assign{\Obj <}{N} = \Setabs{\tuple{n, m}}{n \in \Nat, m \in
  \Nat, n < m}$
\end{enumerate}
De zo vastgelegde structuur~$\Struct N$ voor $\Lang L_A$ heet het
\emph{standaardmodel van de rekenkunde}. Het domein bestaat dan uit de
natuurlijke getallen. De niet-logische constanten van~$\Lang L_A$
krijgen dan precies hun gebruikelijke rekenbetekenis.

Dit is lang niet de enige mogelijke !!{structure} voor~$\Lang L_A$.
We kunnen bijvoorbeeld de verzameling~$\Int$ van gehele getallen
gebruiken in plaats van~$\Nat$. Dan passen we de betekenis van $\Obj
0$, $\Obj \prime$, $\Obj +$, $\Obj \times$ en $\Obj <$ daarbij aan.
Maar we kunnen voor~$\Lang L_A$ zelfs structuren kiezen die helemaal niets met
getallen te maken hebben.
\end{ex}

\begin{ex}
Voor een structuur~$\Struct M$ voor de taal~$\Lang L_Z$ van de
verzamelingenleer zijn maar twee dingen nodig: een verzameling en één
tweeplaatsige relatie. Strikt genomen kan daarom de verzameling mensen
met de relatie ``$x$ is ouder dan $y$'' als !!a{structure} voor $\Lang
L_Z$ dienen. Ook $\Nat$ met $n \ge m$ voor $n, m \in \Nat$ kan zo'n
structuur zijn.

Bij een bijzonder interessante !!{structure} voor $\Lang L_Z$ zijn de
!!{element}s van het domein zelf verzamelingen. Bovendien betekent
$\Obj \in$ daar echt ``$x$ is !!a{element} van~$y$''. Dit is de
!!{structure}~$\Struct{{HF}}$ van \emph{erfelijk eindige
verzamelingen}: eindige verzamelingen waarvan de elementen, de elementen
daarvan enzovoort ook eindig zijn.
\begin{enumerate}
\item $\Domain{{{HF}}} = \emptyset \cup \Pow{\emptyset} \cup
  \Pow{\Pow{\emptyset}} \cup \Pow{\Pow{\Pow{\emptyset}}} \cup \dots$;
\item $\Assign{\Obj \in}{{{HF}}} = \Setabs{\tuple{x, y}}{x, y \in
  \Domain{{{HF}}}, x \in y}$.
\end{enumerate}
\end{ex}

\begin{digress}
Onze afspraken over wat als !!a{structure} telt, veranderen welke
zinnen in onze logica geldig zijn. We staan geen leeg domein toe, dus
ieder domein bevat minstens één object. Bij de gebruikelijke regels
voor waarheid van gekwantificeerde zinnen is
$\lexists[x][(!A(x) \lor \lnot !A(x))]$ daarom in iedere toegestane
structuur waar. De zin is dus geldig: hij is een logische waarheid.
Ook moet iedere !!{constant} een object in het domein aanwijzen.
Daardoor is de stap van $!A(a)$ naar $\lexists[x][!A(x)]$ een geldige
manier van redeneren. Deze stap heet existentiële generalisatie.
Zouden we toestaan dat een naam iets buiten het domein aanwijst, of
helemaal niets aanwijst, dan gaan we in de richting van
\emph{vrije logica}. Daar is voor dezelfde stap een extra premisse
nodig. Naast $!A(a)$ moeten we ook hebben dat
$\lexists[x][\eq[x][a]]$. Pas uit die twee premissen mogen we dan
$\lexists[x][!A(x)]$ afleiden.
\end{digress}

\end{document}
